\section{Experimental Design}
\label{sec:experimental-design}

Our experimental design answers the following research questions about \ourmethod:
\begin{itemize}[noitemsep,topsep=1pt,parsep=2pt,partopsep=0pt,leftmargin=2.5em]
    \item[\textbf{RQ1}] Can \ourmethod bring benefits for domains where ground-truth verification exists but is expensive?
    \item[\textbf{RQ2}] Can learned evaluation improve agents in domains with no objective evaluation?

    \item[\textbf{RQ3}] Do evaluator replacements guide the search toward better task agents over time?
    \item[\textbf{RQ4}] Can co-evolution improve the evaluators themselves?
 
\end{itemize}



\noindent\textbf{Domains and ground-truth anchors.}~~Each domain pairs a generator role with a learned evaluator role and a \emph{ground-truth anchor} (\cref{fig:domain-utilities}). Paper writing is paired with paper review,  anchored to \apres accept/reject decisions~\citep{zhao2026apres}; we construct a matching writer dataset of \apres titles and abstracts to align train, validation, and test distributions across roles. Proof writing is paired with  proof grading, anchored to \imogradbench human grades~\citep{luong2025towards}, where a proof is accepted only when the epoch's frozen grader awards full credit ($7$ of $7$). Coding uses two anchors: the coding agent is anchored to executable \polyglot tests~\citep{gauthier2024polyglot}, while a co-evolved code reviewer is anchored to \crave~\citep{zhang2025crave}, a dataset of accepted/rejected pull requests, and scores each \polyglot patch at generation time.  Across domains, every role agent is initialized from a minimal template~(\cref{app:prompts}). The long runtime of \swebench~\citep{jimenez2024swebench} precluded its inclusion in this version of our work.

\noindent\textbf{Specialists and generalists.}~~Unlike prior methods, \ourmethod optimizes agents against multiple utility functions simultaneously. This raises the question of which utilities to consider when selecting a best-belief agent. We define a \emph{specialist} as the agent with the highest best-belief on one target utility, and a \emph{generalist} as the agent with the highest average best-belief across active utilities. Reporting both shows whether target-task performance comes from task-specific or joint optimization.

\noindent\textbf{Baselines.}~~Every domain uses a learned baseline, \hgmh, which replaces the \dgm search algorithm of \hyperagents~\citep{zhang2026hyperagents} with the more sample-efficient \hgm, while keeping evaluators frozen.  In each domain, \hgmh uses a fixed external evaluator from prior work. We use the \sakana prompt reviewer~\citep{lu2024ai} for paper  writing, \proofautograder~\citep{luong2025towards} for proof grading, and the \imocode prover~\citep{huang2025winninggoldimo2025}. We additionally evaluate the best published \hyperagents~(\dgmh) patches from \citet{zhang2026hyperagents} for paper reviewing and grading as a proxy for \dgm-level performance, though trained on different foundation models.  \citet{zhang2026hyperagents} did not publish the best \polyglot agent.


\noindent\textbf{Principled search interventions.}~~The co-evolutionary framework permits principled modifications of the search objective at epoch boundaries while maintaining epoch-local convergence guarantees. In the paper domain, we exploit this to correct for LLM self-preference bias~(\cref{sec:related}): after the first replacement, papers accepted by the displaced reviewer form an adversarial pool, and the subsequent epoch additionally rewards rejecting these writer-generated papers while maintaining accuracy on the human \apres data. Evaluator replacement itself stays anchored to \apres~(\cref{app:exp-details}). Additional ablations that isolate each part of \ourmethod are reported in \cref{app:nemotron-ablations}. We selected \gptlow (low) for the main text experiments, as it provides a desirable cost-intelligence trade-off~(\cref{app:exp-details}). 

\noindent\textbf{Costs and reporting.}~~We call the agent with the highest $\epsilon$-best-belief score reached during search the \emph{best-belief agent}. This score is comparable across epochs only when it rests on a fixed ground-truth anchor, so a \emph{global} best-belief winner exists only for anchored roles: the learned evaluators and the verifiable coding agent. The paper writer and proof prover have no anchor and admit only \emph{epoch-local} winners, each the best agent under its epoch's frozen evaluator; we select these per epoch and score them post-hoc against fixed external judges~(\cref{tab:paper-writer-cross-reviewer,tab:imo-proof-cross-grader}). Search budgets are matched across runs, so cost differences reflect search dynamics \textbf{rather than unequal allocation}, and we compare runs along three axes: best agent at matched compute, cost to matched quality, and best overall.

Because equal evaluation budgets can hide large compute differences when evaluation costs vary, we report \emph{blended tokens}: input plus output, with output weighted at $5\times$ input cost~\citep{openai2026pricing,anthropic2026pricing}, counting both generation and evaluation calls. Following \citet{wang2025huxley}, a best-belief agent's cost is the budget consumed when it first attains its highest $\epsilon$-best-belief score, \textbf{not when the epoch or run finishes}. Uncertainty is reported with $95\%$ central Beta (Jeffreys) intervals; further details are in \cref{app:exp-details}.

\subsection{Visualization Guide}
\label{exp_design:reading_results_figures}
All result figures~(\cref{fig:intro-token-efficiency,fig:prover-grader-results,fig:paper-writer-reviewer-results}) share one scheme, comparing several \emph{arms}: search configurations, such as the \ourmethod generalist, specialist, and \hgmh baseline. Each figure pairs a \emph{ground-truth panel}, which scores arms against the fixed anchor and so permits direct comparison, with a \emph{search-trajectory panel} tracking best-belief utility as search consumes tokens. A \emph{circle} marks a task agent (writer, prover, or coder) and a \emph{down-triangle} a learned evaluator (reviewer or grader), color identifying the arm. The \emph{crowned heart} marks the single global best-belief winner shown in the ground-truth panel.

The search-trajectory panel instead plots each task agent's best-belief utility against its \emph{own} evaluator, so heights are not comparable across arms; the cross-arm winner is read from the ground-truth panel and tables. Markers along a trajectory are that arm's epoch-local best-belief picks, later scored on held-out test data or fixed evaluators. A \emph{crown on a dashed rule} marks an evaluator replacement, and alternating \emph{shaded bands} the epochs. At each replacement the best-belief may drop, as selective erasure discards records that depended on the displaced evaluator, then re-climbs under the new, typically stricter criterion. A best-belief of $0$ marks the start of an epoch, before any record survives the new evaluator, or an agent that fails a task outright, such as code that does not compile. A trajectory ends where its best-belief agent stops changing, not because search halts (total budget is matched) but because no later agent overtakes it; a faint dotted line shows the continued trajectory. Each panel carries two legends: a main one inside for the searched agents, and a strip above for baselines and paired-role arms. Arrows mark each metric's direction, $\boldsymbol{\uparrow}$ higher-is-better and $\boldsymbol{\downarrow}$ lower-is-better; calibration axes with no preferred direction, such as acceptance rate and rank correlations, carry none.